\chapter{Deep Audit: Yul Helper Functions}
\label{ch:deep-audit}

This chapter performs a function-by-function deep read of the generated
Yul helpers, logging for each one: (i)~the excerpt, (ii)~its Rust
counterpart, (iii)~every deviation or slight difference in the
translation, (iv)~any missing check, and (v)~any gas / contract-space /
robustness / auditability improvement. Findings are graded
\textcolor{tracegreen}{\textbf{[OK]}} (faithful, no action),
\textcolor{mathpurple}{\textbf{[NOTE]}} (benign deviation, documented),
\textcolor{codegenblue}{\textbf{[IMPROVE]}} (concrete optimisation or
hardening opportunity), and
\textcolor{rustorange}{\textbf{[RISK]}} (latent correctness/robustness
concern).

% ---------------------------------------------------------------------
\section{\texttt{AssemblyHelpers.yul}}
\label{sec:asm-helpers}

This partial defines the scalar, transcript, and EC primitives that every
other block calls. It is the highest-leverage file in the generated
contract: a bug here breaks all checks at once.

% ---------------------------------------------------------------------
\subsection{\texttt{scalar\_inv}}
\label{sec:scalar-inv}

\begin{lstlisting}[style=yul]
function scalar_inv(x) -> inv {
    if iszero(x) { revert(0, 0) }
    let p := <scalar_inv_scratch_mptr>
    mstore(add(p, BASE_LEN_OFF), 32)   // base len
    mstore(add(p, EXP_LEN_OFF), 32)    // exp len
    mstore(add(p, MOD_LEN_OFF), 32)    // mod len
    mstore(add(p, BASE_OFF), x)
    mstore(add(p, EXP_OFF),  sub(FR_MODULUS, 2))
    mstore(add(p, MOD_OFF),  FR_MODULUS)
    if iszero(staticcall(gas(), 0x05, p, 0xc0, p, 0x20)) { revert(0, 0) }
    if iszero(eq(returndatasize(), 0x20)) { revert(0, 0) }
    inv := mload(p)
}
\end{lstlisting}

\textbf{Rust counterpart.} \rust{Field::invert} (via \rust{ff}'s
\texttt{Invert} trait), which returns a \rust{CtOption}: zero maps to
\texttt{None}, non-zero to the inverse.

\begin{itemize}
  \item \textcolor{mathpurple}{\textbf{[NOTE] Zero handling differs.}}
        Rust returns \texttt{None} on zero; the Yul \emph{reverts}.
        For a verifier this is equivalent-or-stronger (fail-closed), but
        it is a behavioural difference: a caller that in Rust would
        branch on \texttt{CtOption} instead hard-reverts in Yul. No
        accepted proof changes.
  \item \textcolor{tracegreen}{\textbf{[OK] Precompile frame.}} The
        EIP-198 layout (base/exp/mod lengths then operands) and the
        \texttt{r-2} exponent are correct for BLS12-381 \texttt{Fr}.
  \item \textcolor{tracegreen}{\textbf{[OK] Return-size guard.}} The
        explicit \texttt{returndatasize()==32} check prevents reading
        stale memory if the precompile returns short.
  \item \textcolor{codegenblue}{\textbf{[IMPROVE] Dead-code check.}}
        \texttt{scalar\_inv} duplicates the single-element fast path of
        \texttt{batch\_invert} (\S\ref{sec:batch-invert}) but with a
        \emph{different} zero policy (revert vs.\ return~0). It
        \emph{is} emitted: B64 calls it five times in the
        $f_\text{eval}$ Horner (L3136, L3180, L3216, L3245, L3265), so
        the inconsistent zero policy is live. Batching those five
        inverses into one \texttt{batch\_invert} (\S\ref{sec:roadmap-p3})
        would make \texttt{scalar\_inv} dead and remove it; otherwise
        make the two zero policies identical. \vstatus{OPEN}
\end{itemize}

% ---------------------------------------------------------------------
\subsection{Transcript helpers: \texttt{transcript\_init}, \texttt{common\_word}}
\label{sec:transcript-helpers}

\begin{lstlisting}[style=yul]
function transcript_init() -> buf_len {
    buf_len := TRANSCRIPT_MPTR
}
function common_word(buf_len, word) -> ret {
    mstore(buf_len, word)
    ret := add(buf_len, 32)
}
\end{lstlisting}

\textbf{Rust counterpart.} \rust{Keccak256::common\_\allowbreak{}bytes} appends raw
bytes to the running buffer (\S\ref{sec:rust-transcript}).

\begin{itemize}
  \item \textcolor{tracegreen}{\textbf{[OK] Faithful.}} \texttt{common\_word}
        is the 32-byte specialisation of \texttt{common\_bytes}; the
        cursor arithmetic matches the Rust \texttt{Vec} append exactly.
  \item \textcolor{rustorange}{\textbf{[RISK] No bounds guard.}} The Yul
        helper trusts the codegen-planned buffer size; there is no
        runtime check that \texttt{buf\_len} stays inside the transcript
        region. The Rust \texttt{Vec} grows dynamically and cannot
        overflow. This is safe only if the codegen sizes the buffer
        (\texttt{TranscriptBufferLayout}) correctly, and it once did
        not: before \texttt{6b3d2096} the pre-first-squeeze run counted
        only phase-0 advice, so a valid shape with advice in an early
        phase and the first challenge in a later one overran
        \texttt{VK\_MPTR} and overwrote loaded VK words with proof bytes
        (fix: \texttt{src/\allowbreak{}lowering/\allowbreak{}abi/\allowbreak{}proof.rs} L318--347; B64 has one
        user phase and is unaffected). \emph{Defense-in-depth:}
        \texttt{if gt(buf\_\allowbreak{}len, VK\_\allowbreak{}MPTR) \{ revert(0,0) \}} inside
        \texttt{squeeze\_to} and once after the final $\pi$ absorb
        ($<$0.3k gas). A single end-of-absorption check, as first
        suggested here, would miss that bug: every squeeze resets
        \texttt{buf\_len}, so an earlier overrun is gone by the end.
        \vstatus{PARTIAL} sizing fixed, runtime guard still absent
        (\S\ref{sec:roadmap-p2}).
\end{itemize}

% ---------------------------------------------------------------------
\subsection{\texttt{common\_\allowbreak{}uncompressed\_\allowbreak{}g1}}
\label{sec:common-g1}

\begin{lstlisting}[style=yul]
function common_uncompressed_g1(buf_len, cptr) -> ret {
    let x_hi_word := calldataload(cptr)
    let x_lo      := calldataload(add(cptr, 0x20))
    let y_hi_word := calldataload(add(cptr, 0x40))
    let y_lo      := calldataload(add(cptr, 0x60))
    if shr(128, x_hi_word) { revert(0, 0) }   // top 16 bytes must be 0
    if shr(128, y_hi_word) { revert(0, 0) }
    let x_hi := and(x_hi_word, 0xffff...ffff)  // low 128 bits
    let y_hi := and(y_hi_word, 0xffff...ffff)
    if iszero(or(lt(x_hi, BLS_P_HI),
        and(eq(x_hi, BLS_P_HI), iszero(gt(x_lo, BLS_P_MINUS_ONE_LO))))) {
        revert(0, 0)
    }
    if iszero(or(lt(y_hi, BLS_P_HI),
        and(eq(y_hi, BLS_P_HI), iszero(gt(y_lo, BLS_P_MINUS_ONE_LO))))) {
        revert(0, 0)
    }
    calldatacopy(buf_len, cptr, 0x80)
    ret := add(buf_len, 0x80)
}
\end{lstlisting}

\textbf{Rust counterpart.} \rust{Hashable<Keccak256> for G1Projective}
via \texttt{to\_input} (\texttt{transcript/\allowbreak{}implementors.rs}), which
emits the 128-byte EIP-2537 padded uncompressed form.

\begin{itemize}
  \item \textcolor{tracegreen}{\textbf{[OK] Canonicality check is correct.}}
        The two-level comparison
        $x < p \iff x_{hi} < p_{hi} \lor (x_{hi}=p_{hi} \land x_{lo} < p_{lo})$
        is implemented as
        \texttt{or(lt(x\_\allowbreak{}hi,P\_\allowbreak{}HI), and(eq(x\_\allowbreak{}hi,P\_\allowbreak{}HI), iszero(gt(x\_\allowbreak{}lo,P\_\allowbreak{}MINUS\_\allowbreak{}ONE\_\allowbreak{}LO))))}.
        Since \texttt{P\_MINUS\_ONE\_LO} $= p_{lo}-1$, the term
        \texttt{iszero(gt(x\_\allowbreak{}lo, P\_\allowbreak{}MINUS\_\allowbreak{}ONE\_\allowbreak{}LO))} is exactly
        $x_{lo} \le p_{lo}-1 \iff x_{lo} < p_{lo}$. We verified the
        boundary $x=p$ is correctly rejected.
  \item \textcolor{tracegreen}{\textbf{[OK] Top-byte rejection.}}
        \texttt{shr(128, x\_\allowbreak{}hi\_\allowbreak{}word)} rejects any non-zero in the 16
        padding bytes, preventing transcript-aliasing from
        non-normalised encodings.
  \item \textcolor{mathpurple}{\textbf{[NOTE] Solidity is \emph{stricter}
        than Rust at absorption.}} The Rust absorbs the raw 128 bytes
        without a canonicality gate; a non-canonical encoding would
        instead cause a downstream transcript mismatch. The Yul fails
        fast. The set of \emph{accepted} proofs is identical (canonical
        only), so this is a strengthening, not a divergence.
  \item \textcolor{rustorange}{\textbf{[RISK] No curve / subgroup check
        here.}} This helper validates the \emph{encoding}, not that the
        point lies on BLS12-381 or in the prime-order subgroup. Safety
        rests entirely on
        \cg{validate\_\allowbreak{}g1\_\allowbreak{}precompile\_\allowbreak{}coverage}
        (\texttt{src/\allowbreak{}lowering/\allowbreak{}kzg/\allowbreak{}mod.rs} L569, run from
        \cg{LoweringPlan::validate\_\allowbreak{}generator\_\allowbreak{}invariants},
        \texttt{plan.rs} L238--251)
        guaranteeing every absorbed point later reaches a
        \textbf{G1MSM or pairing} (which \emph{do} subgroup-check per
        EIP-2537). \textbf{G1ADD alone does NOT subgroup-check}, so a
        future codegen change that routes an absorbed point only through
        \texttt{G1ADD} would silently drop the subgroup check.
        \emph{Defense-in-depth:} add an explicit on-curve check
        ($y^2 \equiv x^3+4$ via \texttt{mulmod}) at absorption, or
        assert in the codegen test that no absorbed point's only
        consumer is \texttt{G1ADD}. The second option already holds:
        the coverage list (\cg{precompile\_\allowbreak{}validated\_\allowbreak{}g1s}) admits only
        G1MSM and pairing consumers, so a G1ADD-only point fails the
        build. \vstatus{ADDRESSED} invariant enforced; off-curve
        mutation test \texttt{src/test.rs} L2408. The list is derived
        from the protocol plan, not parsed from the emitted Yul.
\end{itemize}

% ---------------------------------------------------------------------
\subsection{\texttt{squeeze\_to}}
\label{sec:squeeze-to}

\begin{lstlisting}[style=yul]
function squeeze_to(buf_len, mptr) -> ret {
    let h0 := keccak256(TRANSCRIPT_MPTR, sub(buf_len, TRANSCRIPT_MPTR))
    mstore(TRANSCRIPT_MPTR, h0)        // reseed buffer with digest
    let r := FR_MODULUS
    mstore(mptr, mod(h0, r))          // sample Fq = uint256(digest) mod r
    ret := add(TRANSCRIPT_MPTR, 32)
}
\end{lstlisting}

\textbf{Rust counterpart.} \rust{Keccak256::squeeze}
(\texttt{implementors.rs} L186): digest, reseed, then
\rust{sample\_\allowbreak{}keccak\_\allowbreak{}digest\_\allowbreak{}be\_\allowbreak{}mod\_\allowbreak{}r} (\texttt{implementors.rs}
L19), a big-endian digest mod $r$ (\S\ref{sec:rust-transcript},
Def.~\ref{def:squeeze}).

\begin{itemize}
  \item \textcolor{tracegreen}{\textbf{[OK] Faithful.}} Digest, reseed,
        and reduction match the Rust exactly.
  \item \textcolor{mathpurple}{\textbf{[NOTE] Modulo bias.}}
        $\mod(h_0, r)$ with $h_0 \in [0,2^{256})$ and
        $r \approx 2^{254.86}$ gives a non-uniform sample:
        $2^{256}=2r+t$ with $t\approx0.21r$, so residues below $t$ have
        three preimages and the rest two ($1.5\times$ more likely). The
        Rust does the same. \vstatus{CORRECTED} The bias is \emph{not}
        $\le 2^{-128}$: the statistical distance from uniform is
        $\approx0.075$ and the maximum probability is $\approx1.36/r$.
        It is still not exploitable: min-entropy is $\approx254.4$ bits,
        so any Schwartz--Zippel bound grows by at most $1.36\times$.
  \item \textcolor{codegenblue}{\textbf{[IMPROVE] Reseed is a full
        32-byte overwrite.}} Correct, but the buffer length after reseed
        is always \texttt{TRANSCRIPT\_\allowbreak{}MPTR+32}; the codegen could track
        this as a constant rather than recomputing
        \texttt{add(TRANSCRIPT\_\allowbreak{}MPTR,32)} at every squeeze. Negligible
        gas; purely an audit-clarity note.
\end{itemize}

% ---------------------------------------------------------------------
\subsection{\texttt{batch\_invert}}
\label{sec:batch-invert}

\begin{lstlisting}[style=yul]
function batch_invert(success, mptr_start, mptr_end, scratch_mptr, r) -> ret {
    ret := success
    if iszero(ret) { leave }
    if lt(mptr_end, mptr_start) { ret := 0; leave }
    let count_bytes := sub(mptr_end, mptr_start)
    if iszero(count_bytes) { leave }                 // empty batch OK
    if eq(count_bytes, 0x20) {                      // single-element fast path
        let x := mload(mptr_start)
        if iszero(x) { ret := 0; leave }
        ... modexp(x, r-2, r) in place ...
        ret := and(ret, eq(returndatasize(), 0x20))
        if ret { mstore(mptr_start, mload(single_scratch)) }
        leave
    }
    // Forward pass: prefix products (n-1 of them) into scratch; gp = total.
    let gp_mptr := scratch_mptr
    let gp := mload(mptr_start)
    let mptr := add(mptr_start, 0x20)
    for {} lt(mptr, sub(mptr_end, 0x20)) {} {
        gp := mulmod(gp, mload(mptr), r)
        mstore(gp_mptr, gp)
        mptr := add(mptr, 0x20); gp_mptr := add(gp_mptr, 0x20)
    }
    gp := mulmod(gp, mload(mptr), r)
    if iszero(gp) { ret := 0; leave }               // some denominator was 0
    ... modexp(gp, r-2, r) -> all_inv ...
    // Backward pass: recover each inverse from all_inv and prefix products.
    ...
}
\end{lstlisting}

\textbf{Rust counterpart.} \rust{ff::batch\_invert} / the verifier's
own Montgomery batch used for Lagrange denominators and the
$1/(\text{eval}+\beta s+\gamma)$ family.

\begin{itemize}
  \item \textcolor{tracegreen}{\textbf{[OK] Montgomery trick is correct.}}
        We hand-traced $n=2$ and $n=3$: the forward pass stores
        $n-1$ prefix products, the single modexp inverts the total, and
        the backward pass recovers each $a_i^{-1}$ with two
        \texttt{mulmod}s. The special-cased first/second elements at the
        tail are correct.
  \item \textcolor{tracegreen}{\textbf{[OK] Single-element fast path.}}
        Avoids the prefix scratch entirely for one denominator --- a
        real gas win (saves the forward/backward loop and the scratch
        writes).
  \item \textcolor{tracegreen}{\textbf{[OK] Zero-total detection.}} A
        zero total product (any zero denominator) returns
        \texttt{ret=0} rather than producing a bogus inverse. Matches
        the Rust \texttt{CtOption} semantics.
  \item \textcolor{rustorange}{\textbf{[RISK] Memory left dirty on
        failure.}} On a zero total the function \texttt{leave}s with
        only the scratch written; on a failed or short modexp it does
        \emph{not} leave, and the backward pass overwrites the whole
        target range with garbage (template
        \texttt{AssemblyHelpers.yul} L189--208). \vstatus{CORRECTED}
        The caller does \emph{not} gate before reading: the only call
        site (\texttt{Lagrange.yul} L34) reads the inverted range at
        L38--82 and gates at L89 (B64 L1350--1394, gate L1397). This is
        safe because \texttt{success} is sticky and every
        \texttt{success}$=0$ path reverts before \texttt{return}, so
        garbage never reaches an accept decision. \emph{Hardening:}
        document that ``check-before-return'' invariant at the function
        header (L124--126 only says it returns a boolean). A
        ``gate-before-read'' codegen assertion, as first proposed here,
        would reject the current Lagrange block. \vstatus{OPEN}
  \item \textcolor{codegenblue}{\textbf{[IMPROVE] \texttt{gp\_mptr}
        underflow is benign but ugly.}} In the $n=2$ case the backward
        pass sets \texttt{gp\_\allowbreak{}mptr := sub(scratch, 0x20)} (below the
        scratch region) before a loop that never executes. It is
        harmless (the pointer is only used inside the loop) but reads as
        a bug to an auditor. \emph{Rewrite:} restructure the backward
        loop so the pointer is only decremented when the loop body runs.
\end{itemize}

% ---------------------------------------------------------------------
\subsection{\texttt{ec\_pairing}}
\label{sec:ec-pairing}

\begin{lstlisting}[style=yul]
function ec_pairing(success, lhs_mptr, rhs_mptr) -> ret {
    ret := success
    if iszero(ret) { leave }
    let scratch := <final_pairing_scratch_mptr>
    mcopy(scratch,            lhs_mptr,           0x80)
    mcopy(add(scratch, 0x80),  G2_BASE_MPTR,     0x100)
    mcopy(add(scratch, 0x180), rhs_mptr,         0x80)
    mcopy(add(scratch, 0x200), NEG_S_G2_BASE_MPTR, 0x100)
    ret := staticcall(gas(), 0x0f, scratch, 0x300, scratch, 0x20)
    ret := and(ret, eq(returndatasize(), 0x20))
    ret := and(ret, mload(scratch))
    if iszero(ret) { revert(0, 0) }
    ret := 1
}
\end{lstlisting}

\textbf{Rust counterpart.} \rust{DualMSM::check}
(\texttt{msm.rs} L284--300):
\texttt{multi\_\allowbreak{}miller\_\allowbreak{}loop} of $(\text{left}, s\_g2)$ and
$(\text{right}, -g_2)$, final exponentiation, compare to identity
(\S\ref{sec:rust-dualmsm}, Eq.~\eqref{eq:final-pairing-dual}).

\begin{itemize}
  \item \textcolor{tracegreen}{\textbf{[OK] Pairing identity correct.}}
        Checks $e(\text{lhs},[1]_2)\cdot e(\text{rhs},-[s]_2)=1$,
        i.e.\ $e(\text{lhs},[1]_2)=e(\text{rhs},[s]_2)$. The
        \texttt{NEG\_S\_G2\_BASE} is the pre-negated $-[s]_2$ from the
        VK.
  \item \textcolor{mathpurple}{\textbf{[NOTE] Argument-order trap.}}
        \texttt{lhs} is the pairing \emph{first} argument
        ($e(\text{lhs},[1]_2)$) but in the dual-MSM accumulator
        ``left'' historically means $\pi$. The generated comment warns
        that the caller passes \texttt{(success, PAIRING\_\allowbreak{}RHS,
        PAIRING\_\allowbreak{}LHS)} --- swapped relative to the accumulator naming.
        This is a documented, deliberate inversion; it is the single
        most confusing line in the contract and a prime audit target.
  \item \textcolor{tracegreen}{\textbf{[OK] MCOPY layout.}} Cancun
        \texttt{mcopy} replaces the old mstore chains; the comment
        claims $\sim$500 gas saved per call. The 4-copy layout
        (G1|G2|G1|G2) matches the EIP-2537 pairing input ABI.
  \item \textcolor{tracegreen}{\textbf{[OK] G2 subgroup safety.}} The
        pairing precompile checks G2 subgroup membership, so a malicious
        VK \texttt{G2\_BASE} would fail the pairing (rejecting all
        proofs --- a liveness, not soundness, issue). The VK is pinned
        by \texttt{EXPECTED\_\allowbreak{}VK\_\allowbreak{}CODEHASH}, so the G2 is trusted.
  \item \textcolor{mathpurple}{\textbf{[NOTE] Reverts, does not return
        0.}} Unlike \texttt{batch\_invert} (returns bool),
        \texttt{ec\_pairing} \texttt{revert}s on failure (line
        \texttt{if iszero(ret) revert}). Both fail closed, but the two
        helpers use inconsistent failure conventions. For uniform
        \texttt{success}-plumbing and auditability, pick one convention
        (return-bool and let the boundary revert) and apply it to both.
\end{itemize}

% ---------------------------------------------------------------------
\subsection{Summary of \texttt{AssemblyHelpers.yul} findings}

\begin{center}
\footnotesize
\begin{tabular}{@{}p{0.20\linewidth}p{0.09\linewidth}p{0.45\linewidth}p{0.18\linewidth}@{}}
\toprule
\textbf{Function} & \textbf{Grade} & \textbf{Key finding} & \textbf{Status} \\
\midrule
\texttt{scalar\_inv} & \textcolor{codegenblue}{IMPROVE} &
  Duplicates \texttt{batch\_invert} fast path with a different zero
  policy (revert vs.\ return 0). Not dead code: five call sites in B64. &
  \vstatus{OPEN} \\
\texttt{transcript\_init} / \texttt{common\_word} & \textcolor{rustorange}{RISK} &
  Faithful arithmetic, but the cursor is unbounded and the codegen once
  under-sized the buffer; add a per-squeeze bound check. &
  \vstatus{PARTIAL}\newline sizing fixed \texttt{6b3d2096} \\
\texttt{common\_\allowbreak{}uncompressed\_\allowbreak{}g1} & \textcolor{rustorange}{RISK} &
  Encoding canonicality correct; \textbf{no on-curve/subgroup check}
  --- load-bearing on MSM/pairing coverage invariant. &
  \vstatus{ADDRESSED}\newline invariant enforced at build \\
\texttt{squeeze\_to} & \textcolor{tracegreen}{OK} &
  Faithful; modulo bias is real (SD $\approx0.075$) but not exploitable
  and identical to Rust. &
  \vstatus{CORRECTED}\newline was ``$\le2^{-128}$'' \\
\texttt{batch\_invert} & \textcolor{rustorange}{RISK} &
  Correct Montgomery trick; leaves memory dirty on failure. The caller
  reads before its gate; the load-bearing invariant is
  check-before-return. &
  \vstatus{OPEN}\newline premise corrected \\
\texttt{ec\_pairing} & \textcolor{mathpurple}{NOTE} &
  Correct identity; argument-order trap documented; failure convention
  inconsistent with \texttt{batch\_invert}. &
  \vstatus{OPEN} \\
\bottomrule
\end{tabular}
\end{center}

% =====================================================================
\section{\texttt{Lagrange.yul}}
\label{sec:lagrange-audit}

Computes $x^n$, the Lagrange basis $\{L_i(x)\}$, the blinding sum
$l_{\mathsf{blind}}$, and the instance polynomial evaluation, using one
batch inversion.

\begin{lstlisting}[style=yul]
// x^n by repeated squaring, n = 2^k
let x_n := x
for { let idx := 0 } lt(idx, k) { idx := add(idx, 1) } {
    x_n := mulmod(x_n, x_n, r)
}
// Denominators (x - omega_i) and (x^n - 1) batch-inverted in one modexp
for { let pow_of_omega := mload(OMEGA_INV_TO_L_MPTR) }
    lt(mptr, mptr_end) { mptr := add(mptr, 0x20) } {
    mstore(mptr, addmod(x, sub(r, pow_of_omega), r))
    pow_of_omega := mulmod(pow_of_omega, omega, r)
}
mstore(mptr_end, addmod(x_n, sub(r, 1), r))
success := batch_invert(success, X_N_MPTR, add(mptr_end, 0x20), BATCH_INV_SCRATCH_MPTR, r)
// L_i(x) = (x^n - 1) * n^-1 * omega_i / (x - omega_i)
let l_i_common := mulmod(x_n_minus_1, mload(N_INV_MPTR), r)
...
\end{lstlisting}

\textbf{Rust counterpart.} \rust{lagrange\_\allowbreak{}interpolate} /
\rust{lagrange\_\allowbreak{}base\_\allowbreak{}eval} and the instance-eval sum in
\texttt{verifier.rs} (Eq.~\eqref{eq:lagrange}, \eqref{eq:instance-eval}).

\begin{itemize}
  \item \textcolor{tracegreen}{\textbf{[OK] $x^n$ by squaring.}} $k$
        squarings give $x^{2^k}=x^n$; matches the Rust domain exponent.
  \item \textcolor{tracegreen}{\textbf{[OK] Batched denominators.}} All
        $(x-\omega_i)$ plus $(x^n-1)$ are inverted in a single
        \texttt{batch\_invert} (one modexp), matching the Rust's batched
        Lagrange denominators.
  \item \textcolor{tracegreen}{\textbf{[OK] $L_i$ formula.}}
        $L_i(x)=(x^n-1)\,n^{-1}\,\omega_i/(x-\omega_i)$ is exactly
        \eqref{eq:lagrange}.
  \item \textcolor{mathpurple}{\textbf{[NOTE] \texttt{num\_instances==0}
        slot.}} When there are no public instances the code still reserves
        one denominator slot so $L_0$ can be recovered for boundary
        identities. Correct, but the conditional \texttt{mptr\_end} bump
        is easy to misread; it is load-bearing for the
        $l_0/l_{\mathsf{last}}$ reads at the end.
  \item \textcolor{codegenblue}{\textbf{[IMPROVE] \texttt{pow\_of\_omega}
        recomputed twice.}} The first pass walks
        $\omega_{\mathsf{inv}}^L \cdot \omega^i$ to build denominators;
        the second pass re-walks the same sequence to form $L_i$. The
        powers could be cached in the first pass (they are already
        written as the denominator) --- but the second pass needs
        $\omega_i$ not $x-\omega_i$, so a second walk is genuinely
        required. The improvement is only to hoist
        \texttt{l\_i\_common} (already done). Minor.
  \item \textcolor{rustorange}{\textbf{[RISK] Layout-indexed reads.}}
        $l_0$ is read at
        \texttt{add(X\_\allowbreak{}N\_\allowbreak{}MPTR, num\_\allowbreak{}neg\_\allowbreak{}lagranges*32)} and
        $l_{\mathsf{last}}$ at \texttt{X\_N\_MPTR}. These offsets are
        codegen-computed; a mismatch between the Lagrange write layout and
        these read offsets would silently swap $l_0$/$l_{\mathsf{last}}$.
        The differential trace pins them, but the coupling is implicit.
\end{itemize}

% =====================================================================
\section{\texttt{TranscriptProofParser.yul}}
\label{sec:parser-audit}

The full transcript schedule: VK digest, committed identity, instances,
per-phase advice, $\theta$, multiplicities, $\beta,\gamma$, perm $Z$,
lookup helpers/acc, trash, $y$, quotient limbs, $x$, evals, $x_1,x_2$,
$f_{\mathsf{com}}$, $x_3$, $q$-evals, $x_4$, $\pi$.

\begin{lstlisting}[style=yul]
buf_len := common_word(buf_len, mload(VK_DIGEST_MPTR))
// committed_pi = identity -> 128 zero bytes, BEFORE instance count
mstore(buf_len, 0); mstore(add(buf_len,0x20), 0)
mstore(add(buf_len,0x40), 0); mstore(add(buf_len,0x60), 0)
buf_len := add(buf_len, 0x80)
// instance count + values (canonical Fr)
buf_len := common_word(buf_len, <num_instances>)
for ... { let inst_be := calldataload(instance_cptr)
    success := and(success, lt(inst_be, r))
    buf_len := common_word(buf_len, inst_be) }
if iszero(success) { revert(0, 0) }   // fail fast (1e21311c)
...
// final exact-consumption check
if iszero(eq(proof_cptr, NUM_INSTANCE_CPTR)) { revert(0, 0) }
if iszero(success) { revert(0, 0) }
\end{lstlisting}

\textbf{Rust counterpart.} \rust{parse\_trace} +
\rust{verify\_\allowbreak{}algebraic\_\allowbreak{}constraints} transcript schedule
(\S\ref{sec:parse-trace}, \S\ref{sec:rust-transcript}).

\begin{itemize}
  \item \textcolor{tracegreen}{\textbf{[OK] Schedule fidelity.}} The
        absorb/squeeze order matches the Rust exactly; the differential
        trace (ids 1--36) pins every step.
  \item \textcolor{tracegreen}{\textbf{[OK] Committed-identity absorb.}}
        128 zero bytes before the instance count matches the patched
        \texttt{Hashable} identity convention.
  \item \textcolor{tracegreen}{\textbf{[OK] Exact-consumption check.}}
        \texttt{eq(proof\_\allowbreak{}cptr, NUM\_\allowbreak{}INSTANCE\_\allowbreak{}CPTR)} makes any
        under/over-read of the proof layout fail closed --- strong
        defense-in-depth against future layout drift.
  \item \textcolor{tracegreen}{\textbf{[OK] Eval canonicality.}} Every
        proof eval scalar is range-checked \texttt{lt(eval, r)} before
        absorption (lines 319, 382).
  \item \textcolor{rustorange}{\textbf{[RISK] $x_3$ truncation
        (128-bit).}} Under \texttt{truncated-challenges}, $x_3$ is
        masked to 128 bits (line 365). As discussed in
        \S\ref{sec:deviations-summary}, this reduces the
        evaluation-point entropy of the $f$-polynomial check; the
        resulting soundness level is not established here. Faithful to Rust,
        but the single most security-relevant line in the parser.
  \item \textcolor{mathpurple}{\textbf{[NOTE] Deferred \texttt{success}
        for instances only.}} Instance canonicality is accumulated with
        \texttt{success := and(...)} while G1/eval helpers revert
        immediately. \vstatus{ADDRESSED} Since \texttt{1e21311c} the
        parser reverts right after the instance loop (template L68,
        B64 L1050), so a bad instance no longer pays for the whole
        transcript; the final \texttt{if iszero(success) revert} is now
        a backstop.
\end{itemize}

% =====================================================================
\section{\texttt{QuotientAndLinearization.yul}}
\label{sec:quotlin-audit}

External quotient-evaluator call (pinned) plus the linearization scalar
preparation ($x_{\mathsf{split}}$, $1-x^n$).

\begin{lstlisting}[style=yul]
// Pin the external quotient evaluator like the VK
if iszero(and(
    eq(extcodesize(quotientEvaluator), EXPECTED_QUOTIENT_LENGTH),
    eq(extcodehash(quotientEvaluator), EXPECTED_QUOTIENT_CODEHASH_WORD)
)) { revert(0, 0) }
if iszero(staticcall(gas(), quotientEvaluator, frame_base, frame_len, q_out, out_len)) { revert(0,0) }
if iszero(eq(mload(q_out), <magic>)) { revert(0, 0) }
mstore(QUOTIENT_EVAL_MPTR, mload(add(q_out, 0x20)))   // -nu_y(x)
// x_split = x^(n-1), one_minus_x_n = 1 - x^n, one squaring walk
let x_pow_2i := x
let x_pow_2i_minus1 := 1
for { let idx := 0 } lt(idx, k) { idx := add(idx, 1) } {
    x_pow_2i_minus1 := mulmod(mulmod(x_pow_2i_minus1, x_pow_2i_minus1, r), x, r)
    x_pow_2i := mulmod(x_pow_2i, x_pow_2i, r)
}
mstore(QUOTIENT_MPTR, x_pow_2i_minus1)        // x_split
mstore(add(QUOTIENT_MPTR, 0x20), one_minus_x_n)
\end{lstlisting}

\textbf{Rust counterpart.} \rust{compute\_\allowbreak{}linearization\_\allowbreak{}commitment}
(\S\ref{sec:lin-com}); the external path is the split-contract variant
of the same numerator.

\begin{itemize}
  \item \textcolor{tracegreen}{\textbf{[OK] Evaluator pinned by
        codehash.}} The external quotient contract is checked against
        \texttt{EXPECTED\_\allowbreak{}QUOTIENT\_\allowbreak{}CODEHASH} before every call ---
        same defense-in-depth as the VK. A swapped evaluator reverts.
  \item \textcolor{tracegreen}{\textbf{[OK] Magic/version gate.}} The
        return-word-0 magic check catches a wrong-callee or ABI drift.
  \item \textcolor{tracegreen}{\textbf{[OK] $x_{\mathsf{split}}$ walk.}}
        Computing $x^{n-1}$ and $x^n$ in one loop via
        $x^{2^i-1}=(x^{2^{i-1}-1})^2\cdot x$ is correct and saves a
        separate exponentiation.
  \item \textcolor{mathpurple}{\textbf{[NOTE] \texttt{QUOTIENT\_MPTR}
        repurposed as scalar metadata.}} In this path
        \texttt{QUOTIENT\_MPTR} is \emph{not} a G1 point; its two words
        carry $(x_{\mathsf{split}}, 1-x^n)$. The comment warns, but an
        auditor expecting a point here could be misled. Consider a
        distinct slot name (e.g.\ \texttt{LIN\_\allowbreak{}SCALARS\_\allowbreak{}MPTR}) to remove
        the ambiguity.
  \item \textcolor{codegenblue}{\textbf{[IMPROVE] Trace path uses
        \texttt{call} not \texttt{staticcall}.}} When tracing, the
        external quotient call is a \texttt{call} (line 28) rather than
        \texttt{staticcall}. \vstatus{CORRECTED} This is required, not
        an oversight: the trace-mode evaluator emits \texttt{log1}
        trace events (\texttt{Halo2QuotientEvaluator.sol} L120--127),
        and \texttt{LOG} reverts inside a \texttt{STATICCALL} frame.
        Production renders use \texttt{staticcall} (line 30), and B64
        embeds the evaluator, so it makes no external call at all. No
        action.
\end{itemize}

% =====================================================================
\section{\texttt{FinalPairing.yul}}
\label{sec:finalpairing-audit}

Randomized batching of the IVC accumulator equation into the KZG pairing,
then the final \texttt{ec\_pairing}.

\begin{lstlisting}[style=yul]
// alpha = FS(domain || kzg_rhs || kzg_lhs || acc_rhs || acc_lhs)
let acc_pair_alpha := mod(keccak256(batch_ptr, HASH_BYTES), r)
if iszero(acc_pair_alpha) { acc_pair_alpha := 1 }
// PAIRING_RHS += alpha * ACC_RHS   (1-pair G1MSM then G1ADD)
mcopy(batch_ptr, ACC_RHS_MPTR, 0x80)
mstore(add(batch_ptr, 0x80), acc_pair_alpha)
success := staticcall(gas(), G1MSM, batch_ptr, 0xa0, batch_ptr, 0x80)
mcopy(add(batch_ptr, 0x80), PAIRING_RHS_MPTR, 0x80)
success := staticcall(gas(), G1ADD, batch_ptr, 0x100, PAIRING_RHS_MPTR, 0x80)
// ... mirror for LHS ...
success := ec_pairing(success, PAIRING_RHS_MPTR, PAIRING_LHS_MPTR)
\end{lstlisting}

\textbf{Rust counterpart.} The accumulator randomizer in
\texttt{verifier.rs} + \rust{DualMSM::check}
(Eq.~\eqref{eq:acc-batch}, \eqref{eq:final-pairing-dual}).

\begin{itemize}
  \item \textcolor{tracegreen}{\textbf{[OK] Randomized batching prevents
        cancellation.}} Deriving $\alpha$ from the pairing inputs means
        two independently-bad equations cannot be made to cancel: the
        combined equation holds for at most one $\alpha$. This is the
        correct construction (vs.\ naively multiplying equations).
  \item \textcolor{tracegreen}{\textbf{[OK] Zero-$\alpha$ guard.}}
        Mapping $\alpha=0\mapsto 1$ prevents the accumulator equation
        from being silently dropped (which would happen if $\alpha=0$
        zeroed the accumulator contribution).
  \item \textcolor{codegenblue}{\textbf{[IMPROVE] Fold the two 1-pair
        G1MSMs.}} The RHS and LHS accumulator scalings are two separate
        1-pair \texttt{G1MSM} calls. They cannot be merged into one MSM
        (different output points), but each 1-pair MSM could instead be
        a \texttt{G1ADD} of precomputed points if the accumulator
        points were pre-scaled --- not applicable here since $\alpha$ is
        runtime. The current form is near-optimal; the only saving is
        the \texttt{G1ADD} after each MSM, which is unavoidable (a 2-pair
        MSM costs 22{,}776 gas vs.\ 12{,}375 for MSM+ADD). The foldable
        one-pair MSMs are $-vG$ and $x_3\pi$ in the PCS block, which can
        join the fused final MSM (\S\ref{sec:roadmap-p3}, register
        I12).
  \item \textcolor{mathpurple}{\textbf{[NOTE] Argument-order trap
        (again).}} \texttt{ec\_\allowbreak{}pairing(success, PAIRING\_\allowbreak{}RHS,
        PAIRING\_\allowbreak{}LHS)} passes RHS as the pairing's first argument. The
        comment explains the dual-MSM naming inversion. This is the
        second place the LHS/RHS inversion bites (see
        \S\ref{sec:ec-pairing}); a single consistent naming
        (\texttt{PAIRING\_ARG0}/\texttt{ARG1}) would remove the trap.
\end{itemize}

% =====================================================================
\section{\texttt{AccumulatorHelpers.yul}}
\label{sec:acc-audit}

Decodes the shifted radix-$2^{56}$ public-input encoding of the IVC
accumulator's foreign G1 points into EIP-2537 padded form, with
canonical packing checks and mandatory G1MSM validation.

\begin{lstlisting}[style=yul]
function load_acc_coord(src, allow_id, bits, n, base, limbs_per_word) -> ok, hi, lo, is_id {
    ok := check_acc_coord_packing(src, bits, n, limbs_per_word)  // reject unused high bits
    if and(allow_id, iszero(lt(calldataload(src), base))) {
        let adj_hi, adj_lo := load_acc_coord_shifted(src, bits, n, base, limbs_per_word, base)
        is_id := is_bls_p_minus_one(adj_hi, adj_lo)
    }
    hi, lo := load_acc_coord_shifted(src, bits, n, base, limbs_per_word, mul(is_id, base))
    ok := and(ok, or(lt(hi, BLS_P_HI), and(eq(hi, BLS_P_HI), iszero(gt(lo, BLS_P_MINUS_ONE_LO)))))
    if is_bls_p_minus_one(hi, lo) { hi := 0; lo := 0 }        // p-1 -> 0
    else { let next_lo := add(lo, 1); hi := add(hi, lt(next_lo, lo)); lo := next_lo }  // +1 with carry
    ok := and(ok, lt(hi, shl(128, 1)))                        // EIP-2537 48-byte pad
}
\end{lstlisting}

\textbf{Rust counterpart.} \texttt{AssignedForeignPoint<BLS12-381>} /
\texttt{AssignedField::as\_\allowbreak{}public\_\allowbreak{}input} shifted-limb codec.

\begin{itemize}
  \item \textcolor{tracegreen}{\textbf{[OK] Packing canonicality.}}
        \texttt{check\_\allowbreak{}acc\_\allowbreak{}coord\_\allowbreak{}packing} rejects any unused high bits
        in the packed words, so each coordinate has a unique encoding
        before it reaches curve validation.
  \item \textcolor{tracegreen}{\textbf{[OK] Mandatory G1MSM validation.}}
        Every decoded point (even identity, even scalar 0/1) is routed
        through \texttt{G1MSM} so the precompile performs the
        on-curve/subgroup check. The comment explicitly notes skipping
        it would let a malformed non-identity point hide behind scalar 0
        --- correct and important.
  \item \textcolor{tracegreen}{\textbf{[OK] Decoded-zero rejection.}} A
        decoded $(0,0)$ that is \emph{not} the canonical identity
        encoding is rejected (line 208--209), since EIP-2537 reserves
        $(0,0)$ for infinity.
  \item \textcolor{tracegreen}{\textbf{[OK] $+1$ carry into high word.}}
        The \texttt{hi := add(hi, lt(next\_\allowbreak{}lo, lo))} carry logic for the
        coord$-1\to$coord reconstruction is correct across the 256-bit
        boundary.
  \item \textcolor{rustorange}{\textbf{[RISK] Intricate identity-flag
        logic.}} The \texttt{first\_adjust}/\texttt{is\_id} path (only
        $x$ carries the flag; the flag is one radix base added to the
        first packed word) is the most intricate decoding in the
        contract. It is correct for the current codec, but a change to
        the circuit's public-input packing would require a matching
        change here. The \texttt{is\_\allowbreak{}acc\_\allowbreak{}encoded\_\allowbreak{}identity} fast path
        (exact packed-word equality) is the primary guard; the
        coordinate-level flag probe is the fallback. High audit burden
        --- recommend a dedicated unit test enumerating identity vs.\
        non-identity encodings. \vstatus{PARTIAL} The IVC E2E covers
        bad packing and malformed LHS/RHS points
        (\texttt{tests/\allowbreak{}ivc\_\allowbreak{}keccak\_\allowbreak{}solidity.rs} L1417--1447); the
        identity path has only a template-string check
        (\texttt{src/\allowbreak{}lowering/\allowbreak{}tests.rs} L1223).
  \item \textcolor{mathpurple}{\textbf{[NOTE] \texttt{y\_id} discarded.}}
        \texttt{pop(y\_id)} is correct (allow\_id was false) but the
        shared tuple shape is a minor readability wart.
\end{itemize}

\subsection{Accumulator final-pairing batching (verified on \texttt{ivc-keccak})}
\label{sec:acc-pairing-verify}

The poseidon fixture has no accumulator, so the accumulator pairing path
was verified against the \texttt{ivc-keccak} dump, which exercises the
full randomized batching. The accumulator contributes a second pairing
equation that must be checked \emph{together with} the KZG equation
without letting two bad equations cancel.

\paragraph{The batching.} After all four G1 pairing inputs are
materialized, a verifier-local randomizer $\alpha$ is drawn by
Fiat--Shamir over the domain tag and all four points, and the pairing
arguments are folded:
\begin{lstlisting}[style=yul]
// domain || KZG_rhs || KZG_lhs || ACC_rhs || ACC_lhs  (32 + 4*128 = 0x220)
mstore(batch_ptr, "pairing-batch-acc-kzg\0...")
mcopy(batch_ptr+0x20,  PAIRING_RHS_MPTR, 0x80)
mcopy(batch_ptr+0xa0,  PAIRING_LHS_MPTR, 0x80)
mcopy(batch_ptr+0x120, ACC_RHS_MPTR,     0x80)
mcopy(batch_ptr+0x1a0, ACC_LHS_MPTR,     0x80)
let acc_pair_alpha := mod(keccak256(batch_ptr, 0x0220), r)
if iszero(acc_pair_alpha) { acc_pair_alpha := 1 }   // never drop the acc eq
// PAIRING_RHS += alpha * ACC_RHS   (G1MSM then G1ADD)
// PAIRING_LHS += alpha * ACC_LHS   (G1MSM then G1ADD)
\end{lstlisting}

\paragraph{The check.} The final \texttt{ec\_pairing} then verifies
\begin{equation}
e\!\big(\mathsf{KZG_{rhs}} + \alpha\,\mathsf{ACC_{rhs}},\,[1]_2\big)\cdot
e\!\big(\mathsf{KZG_{lhs}} + \alpha\,\mathsf{ACC_{lhs}},\,-[s]_2\big) = 1 .
\label{eq:acc-batch-check}
\end{equation}
By bilinearity this factors as
\begin{equation}
\underbrace{e(\mathsf{KZG_{rhs}},[1]_2)\,e(\mathsf{KZG_{lhs}},-[s]_2)}_{\text{KZG equation}}
\cdot
\underbrace{\big[e(\mathsf{ACC_{rhs}},[1]_2)\,e(\mathsf{ACC_{lhs}},-[s]_2)\big]^{\alpha}}_{\text{accumulator equation}}
= 1 .
\label{eq:acc-batch-factor}
\end{equation}

\begin{itemize}
  \item \textcolor{tracegreen}{\textbf{[OK] Cancellation is prevented.}}
        If both the KZG and accumulator equations hold,
        \eqref{eq:acc-batch-factor} is $1\cdot 1^{\alpha}=1$. If either
        is \emph{false}, the product equals $1$ for \textbf{at most one}
        value of $\alpha\in\mathbb{F}_r$ (a degree-1 polynomial in
        $\alpha$ has $\le 1$ root). Since $\alpha$ is drawn
        \emph{after} all four points are fixed, a prover cannot steer it
        into the cancelling root except with probability $\le 1/r
        \approx 2^{-255}$. This is the correct Schwartz--Zippel
        randomization.
  \item \textcolor{tracegreen}{\textbf{[OK] $\alpha$ drawn after inputs
        are fixed.}} The hash covers the fully materialized
        \texttt{PAIRING\_RHS/LHS} and \texttt{ACC\_RHS/LHS}, so the
        prover commits to the points before $\alpha$ exists --- the
        randomizer cannot be chosen adversarially.
  \item \textcolor{tracegreen}{\textbf{[OK] Zero-$\alpha$ guard.}}
        \texttt{if iszero(alpha) alpha := 1} ensures the accumulator
        equation is never accidentally dropped (an $\alpha=0$ draw would
        collapse \eqref{eq:acc-batch-check} to the KZG equation alone,
        silently ignoring the accumulator).
  \item \textcolor{tracegreen}{\textbf{[OK] Domain-separated $\alpha$.}}
        The \texttt{"pairing-batch-acc-kzg"} tag prevents the $\alpha$
        from being reused across protocols / other transcript contexts.
  \item \textcolor{tracegreen}{\textbf{[OK] Points pre-validated.}}
        \texttt{ACC\_RHS/LHS} were already forced through G1MSM in
        \texttt{validate\_\allowbreak{}public\_\allowbreak{}accumulator} (every decoded point,
        even identity/scalar-0), so the $\alpha$-fold MSMs operate on
        validated points; a malformed accumulator reverts before the
        pairing.
\end{itemize}

\paragraph{Encoding correctness (recap).} The accumulator point
encoding was verified end-to-end: the shifted radix-$2^{56}$ little-endian
limb reconstruction (including the limb straddling the 256-bit
\texttt{hi}/\texttt{lo} split via \texttt{shr(sub(256,shift),limb)}),
the canonical-packing rejection of unused high bits, the
$p-1\!\to\!0$ / \texttt{coord-1}$\to$\texttt{coord} $+1$-with-carry
round-trip, the base-field range check, and the EIP-2537
\texttt{hi}$<2^{128}$ pad check all hold. The identity fast path
(\texttt{is\_\allowbreak{}acc\_\allowbreak{}encoded\_\allowbreak{}identity}, exact packed-word equality) plus
the coordinate-level flag probe cover the point-at-infinity encoding,
and a decoded $(0,0)$ that is not the canonical identity is rejected.

\paragraph{Verdict.} No issues with the accumulator \textbf{encoding}
and no issues with the \textbf{final pairing check} in the generated
Solidity. The only residual concern is the audit-burden of the
identity-flag decode (already flagged [RISK] above), which is a
test-coverage recommendation, not a correctness defect.

% =====================================================================
\section{Deviations summary}
\label{sec:deviations-summary}

Aggregating every deviation logged in this chapter, ranked by security
relevance:

\begin{center}
\footnotesize
\begin{tabular}{@{}p{0.26\linewidth}p{0.09\linewidth}p{0.39\linewidth}p{0.18\linewidth}@{}}
\toprule
\textbf{Deviation} & \textbf{Severity} & \textbf{Assessment} & \textbf{Status} \\
\midrule
128-bit truncated challenges ($x_3$, $x_1^i$, $x_4^i$) &
  \textcolor{rustorange}{HIGH} &
  Soundness level not established; earlier bit-level figures are
  withdrawn. Faithful to Rust; needs a written soundness argument and
  confirmation against the security target. &
  \vstatus{OPEN}\newline policy decision \\
Unbounded transcript cursor &
  \textcolor{rustorange}{MED} &
  Buffer sized by codegen only; a mis-sizing for multi-phase shapes let
  proof bytes overwrite the loaded VK. &
  \vstatus{PARTIAL}\newline fixed in \texttt{6b3d2096}, no runtime
  guard \\
No on-curve/subgroup check at G1 absorption &
  \textcolor{rustorange}{MED} &
  Sound only via MSM/pairing coverage invariant; G1ADD alone does not
  subgroup-check. &
  \vstatus{ADDRESSED}\newline enforced at build \\
\texttt{batch\_invert} leaves memory dirty on failure &
  \textcolor{rustorange}{MED} &
  The caller reads before its gate; the load-bearing invariant is
  check-before-return (sticky \texttt{success}). &
  \vstatus{OPEN}\newline premise corrected \\
Accumulator identity-flag decode intricacy &
  \textcolor{rustorange}{MED} &
  Correct for current codec; high audit burden, needs dedicated tests. &
  \vstatus{PARTIAL}\newline no identity vectors \\
\texttt{QUOTIENT\_MPTR} repurposed as scalar metadata &
  \textcolor{codegenblue}{LOW} &
  Audit-clarity; rename to remove ambiguity. &
  \vstatus{OPEN} \\
Trace path uses \texttt{call} not \texttt{staticcall} &
  \textcolor{codegenblue}{LOW} &
  Not a deviation: trace-mode evaluator emits \texttt{LOG}, which
  \texttt{STATICCALL} forbids; production already uses
  \texttt{staticcall}. &
  \vstatus{CORRECTED} \\
\texttt{scalar\_inv} vs \texttt{batch\_invert} zero-policy mismatch &
  \textcolor{codegenblue}{LOW} &
  Live code (five call sites in B64); inconsistent failure convention. &
  \vstatus{OPEN} \\
PAIRING LHS/RHS naming inversion &
  \textcolor{mathpurple}{LOW} &
  Documented trap; rename to ARG0/ARG1. &
  \vstatus{OPEN} \\
Modulo bias in \texttt{squeeze\_to} &
  \textcolor{mathpurple}{NEGL} &
  Statistical distance $\approx0.075$ (not $\le2^{-128}$); same as
  Rust, not exploitable. &
  \vstatus{CORRECTED} \\
\bottomrule
\end{tabular}
\end{center}

The first four are the ones a security reviewer should focus on. The
remainder are audit-clarity and robustness improvements that do not
change the accepted-proof set.

% =====================================================================
\section{\texttt{QuotientNumeratorBlock.yul} --- the compact quotient VM}
\label{sec:vm-audit}

This is the largest and most intricate generated component (1275 lines).
It reconstructs the $y$-batched numerator $\nu_y(x)$ by interpreting a
bytecode program (\texttt{q\_program}) stored in the \emph{VK payload},
rather than unrolling every identity as Yul. The opcode stream is
VK-pinned; no proof calldata can alter control flow.

\subsection{Stack model}
\begin{lstlisting}[style=yul]
// q_pc: bytecode cursor   q_end: exclusive end
// q_sp: spilled-stack word pointer
// q_top: cached top-of-stack value   q_has_top: liveness flag
let q_pc := q_program_mptr
let q_end := add(q_program_mptr, <program.len>)
let q_sp := <program.stack_mptr>
let q_top := 0
let q_has_top := 0
for { } lt(q_pc, q_end) { } {
    let q_op := byte(0, mload(q_pc))
    q_pc := add(q_pc, 1)
    switch q_op { ... default { revert(0, 0) } }
}
// end-of-program integrity:
if iszero(eq(q_pc, q_end)) { revert(0, 0) }
if q_has_top { revert(0, 0) }
\end{lstlisting}

\begin{itemize}
  \item \textcolor{tracegreen}{\textbf{[OK] Fail-closed default.}} Any
        opcode not in the emitted switch (including the reserved
        \texttt{0x1a}) hits \texttt{default \{ revert \}}.
  \item \textcolor{tracegreen}{\textbf{[OK] End-of-program integrity.}}
        The two post-loop checks --- \texttt{q\_pc == q\_end} (no
        over/under-read) and \texttt{q\_has\_top == 0} (no dangling
        partial expression) --- catch malformed generator output. Added
        in \texttt{33f78d2b} (template L1191--1192; B64 L2569--2570).
  \item \textcolor{tracegreen}{\textbf{[OK] All arithmetic reduced.}}
        Every \texttt{addmod}/\texttt{mulmod} is modulo \texttt{r};
        there is no unchecked \texttt{add}/\texttt{mul} on field values.
  \item \textcolor{tracegreen}{\textbf{[OK] Coefficients are VK-pinned.}}
        Constant-table reads come from \texttt{q\_const\_mptr} (VK
        payload), never from proof calldata --- the bytecode cannot be
        steered by the prover.
\end{itemize}

\subsection{\texttt{FOLD\_MAIN} (0x0a) --- the $y$-batch fold}
\begin{lstlisting}[style=yul]
case 0x0a {
    let q_eval := q_top
    q_has_top := 0
    // qn = qn*y + eval   (reverse Horner, matches Rust)
    mstore(qn, mulmod(mload(qn), y, r))
    mstore(qn, addmod(mload(qn), q_eval, r))
}
\end{lstlisting}
\textbf{Rust counterpart.} The reverse $y$-power fold in
\rust{compute\_\allowbreak{}linearization\_\allowbreak{}commitment}
(\texttt{expressions.iter().rev()}, \texttt{y\_pow *= y}).

\begin{itemize}
  \item \textcolor{tracegreen}{\textbf{[OK] Horner fold correct.}}
        \texttt{qn = qn*y + eval} accumulated over identities in order
        yields $\sum_i y^{m-1-i} e_i = \nu_y(x)$, matching
        \eqref{eq:nu-y}.
  \item \textcolor{codegenblue}{\textbf{[IMPROVE] Redundant
        load/store.}} The fold does \texttt{mstore(qn, mulmod(mload(qn),
        y, r))} then \texttt{mstore(qn, addmod(mload(qn), eval, r))}
        --- two \texttt{mstore} + two \texttt{mload} of the same slot per
        identity. Keeping the accumulator in a register and storing once
        at the end would save one \texttt{mload}+\texttt{mstore}
        ($\sim$6 gas) per identity. Correctness is unaffected.
        \vstatus{OPEN} B64 does not render case \texttt{0x0a}; the same
        pair appears in the native-callback fold snippets
        (\texttt{src/\allowbreak{}lowering/\allowbreak{}quotient.rs} L1018, L1035; e.g.\ B64
        L2005--2006). The saving is small, not ``meaningful''.
\end{itemize}

\subsection{\texttt{FOLD\_SELECTOR} (0x0b) --- sparse selector buckets}
\begin{lstlisting}[style=yul]
case 0x0b {
    let q_sel_idx := shr(16, payload)
    let q_sel_gap := and(payload, 0xffff)
    let q_eval := q_top; q_has_top := 0
    // global accumulator still *y (keeps main identities aligned)
    mstore(qn, mulmod(mload(qn), y, r))
    let q_sel_acc := mload(SELECTOR_ACC_MPTR + 32*sel_idx)
    if q_sel_gap {
        q_sel_acc := mulmod(q_sel_acc, mload(selector_power_mptr + 32*gap), r)
    }
    mstore(q_target_ptr, addmod(q_sel_acc, q_eval, r))
}
\end{lstlisting}
\begin{itemize}
  \item \textcolor{tracegreen}{\textbf{[OK] Global alignment preserved.}}
        The global \texttt{qn} is still multiplied by $y$ even for
        selector identities, so subsequent main identities land at the
        same $y$-powers as Rust's reverse fold.
  \item \textcolor{tracegreen}{\textbf{[OK] Precomputed \texttt{y\textasciicircum gap}.}}
        Using a codegen-precomputed \texttt{y\textasciicircum gap} table
        avoids a runtime \texttt{y\textasciicircum\{-1\}} modexp --- a
        real gas win. The tail update (post-VM) applies the final
        $y^{\text{tail}}$ to align each bucket with the global end.
\end{itemize}

\subsection{Native callbacks (0x19 perm, 0x1f lookup, 0x1b heavy)}
\begin{lstlisting}[style=yul]
case 0x19 {  // native permutation
    q_top := 0; q_has_top := 0          // reset stack at identity boundary
    q_sp := <program.stack_mptr>
    // generated perm fold lines (same snippets as interpreted)
}
\end{lstlisting}
\begin{itemize}
  \item \textcolor{tracegreen}{\textbf{[OK] Identity-boundary reset.}}
        Each native callback clears \texttt{q\_top}/\texttt{q\_has\_top}
        and resets \texttt{q\_sp} before running its generated fold, so
        it never inherits a partial VM stack.
  \item \textcolor{rustorange}{\textbf{[RISK] Empty-stack-boundary
        assumption.}} The reset \emph{silently drops} any live
        \texttt{q\_top}. This is safe only because the codegen emits
        native callbacks exclusively at empty-stack boundaries. If a
        future codegen change emitted one mid-expression, the partial
        value would be lost without error.
        \emph{Defense-in-depth:} add
        \texttt{if q\_\allowbreak{}has\_\allowbreak{}top \{ revert(0,0) \}} at the top of each
        native-callback case, turning the assumption into a checked
        invariant. \vstatus{OPEN} The final-bytecode validator rejects a
        callback on a non-empty stack
        (\texttt{src/\allowbreak{}lowering/\allowbreak{}quotient\_\allowbreak{}numerator/\allowbreak{}vm/\allowbreak{}mod.rs}
        L2241--2246), but the runtime reset is still unconditional
        (template L1057, L1084, L1113; B64 L1934, L2051, L2247).
\end{itemize}

\subsection{ADD/MUL stack underflow}
\begin{lstlisting}[style=yul]
case 0x06 {  // ADD
    q_sp := sub(q_sp, 0x20)
    q_top := addmod(mload(q_sp), q_top, r)
}
\end{lstlisting}
\begin{itemize}
  \item \textcolor{rustorange}{\textbf{[RISK] Unchecked stack
        underflow.}} \texttt{ADD}/\texttt{MUL} assume a spilled operand
        exists (\texttt{q\_\allowbreak{}sp > stack\_\allowbreak{}mptr}). The comment says the
        codegen safety validator guarantees this, but the runtime does
        not check. A validator bug would make \texttt{q\_sp} underflow
        below the stack region and read adjacent memory as an operand.
        \emph{Defense-in-depth:}
        \texttt{if lt(q\_\allowbreak{}sp, add(stack\_\allowbreak{}mptr, 0x20)) \{ revert(0,0) \}}
        before the \texttt{sub}. Costs $\sim$10 gas per ADD/MUL; given
        the validator already prevents it, this is belt-and-suspenders
        for auditability. \vstatus{OPEN} The validator re-decodes the
        final bytes and rejects underflow (\texttt{vm/mod.rs}
        L2213--2216; test \texttt{src/\allowbreak{}lowering/\allowbreak{}tests.rs} L1520); no
        runtime check (template L503--518; B64 renders only ADD,
        L1716--1721).
\end{itemize}

\subsection{Final negation}
\label{sec:vm-final-neg}
\begin{lstlisting}[style=yul]
// -nu_y(x) stored as the linearization expected eval
let linearization_expected_eval := addmod(0, sub(r, mload(qn)), r)
mstore(QUOTIENT_EVAL_MPTR, linearization_expected_eval)
\end{lstlisting}
\begin{itemize}
  \item \textcolor{tracegreen}{\textbf{[OK] Matches Rust.}}
        \texttt{expected\_\allowbreak{}eval -= eval} for \texttt{col\_idx == None}
        produces $-\nu_y(x)$; the $(1-x^n)$ factor is on the commitment
        side (Def.~\ref{def:neg-numer}).
\end{itemize}

% =====================================================================
\section{\texttt{QuotientHelpers.yul}}
\label{sec:qhelpers-audit}

Optional codegen-recognized helper functions, emitted only when the
lowering pass recognizes the corresponding expression shape.

\begin{lstlisting}[style=yul]
function q_pow5(x) -> z {                 // Poseidon S-box x^5
    let x2 := mulmod(x, x, q_r)
    z := mulmod(x, mulmod(x2, x2, q_r), q_r)
}
function q_limb7(x0,...,x6) -> z {       // sum_i base_powers[i]*limb_i
    z := addmod(x0, mulmod(c1, x1, q_r), q_r)
    ...
}
\end{lstlisting}
\begin{itemize}
  \item \textcolor{tracegreen}{\textbf{[OK] \texttt{q\_pow5} correct.}}
        $x^5 = x\cdot(x^2)^2$ with two \texttt{mulmod}s --- the
        minimal addition chain for 5.
  \item \textcolor{mathpurple}{\textbf{[NOTE] Codegen shortcuts, not
        verifier rules.}} These mirror the Rust \emph{Expression tree}
        shapes (Poseidon S-box, foreign-field limb sums). They are a
        rendering optimization; the verifier semantics are unchanged.
        Coefficients come from the VK/program constant table, never
        proof data.
\end{itemize}

% =====================================================================
\section{VM-specific improvement summary}
\label{sec:vm-improvements}

\begin{center}
\footnotesize
\begin{tabular}{@{}p{0.26\linewidth}p{0.10\linewidth}p{0.38\linewidth}p{0.18\linewidth}@{}}
\toprule
\textbf{Item} & \textbf{Grade} & \textbf{Recommendation} & \textbf{Status} \\
\midrule
Native-callback empty-stack assumption & \textcolor{rustorange}{RISK} &
  Add \texttt{if q\_\allowbreak{}has\_\allowbreak{}top revert} at each native-callback entry
  (\texttt{0x19}, \texttt{0x1f}, \texttt{0x1b}). &
  \vstatus{OPEN}\newline build-time check only \\
ADD/MUL unchecked stack underflow & \textcolor{rustorange}{RISK} &
  Add \texttt{if lt(q\_\allowbreak{}sp, stack\_\allowbreak{}base+32) revert} before the pop. &
  \vstatus{OPEN}\newline build-time check only \\
FOLD\_MAIN redundant load/store & \textcolor{codegenblue}{IMPROVE} &
  Keep accumulator in a register; store once at end ($\sim$6 gas per
  folded identity; in B64 via native fold snippets). &
  \vstatus{OPEN} \\
Bytecode control-flow pinned by VK & \textcolor{tracegreen}{OK} &
  Prover cannot alter opcode stream. &
  --- \\
End-of-program integrity checks & \textcolor{tracegreen}{OK} &
  \texttt{q\_pc==q\_end} and \texttt{!q\_has\_top} catch malformed output. &
  \vstatus{ADDRESSED}\newline \texttt{33f78d2b} \\
\bottomrule
\end{tabular}
\end{center}

The two \textcolor{rustorange}{RISK} items are the VM's load-bearing
codegen assumptions (empty-stack native callbacks, non-empty-stack
ADD/MUL). Both are guaranteed at build time by
\cg{validate\_\allowbreak{}quotient\_\allowbreak{}program}, which re-decodes the final bytecode
(\texttt{vm/mod.rs} L2003), rather than checked at runtime; adding the
one-line runtime guards would make them self-evident to an auditor at
negligible cost.

% =====================================================================
\section{Contract templates and deployment-time checks}
\label{sec:contract-templates}

The generated verifier is assembled from three top-level contract
templates plus the Yul partials audited above. This section audits the
contract-level structure and the deployment-time defenses.

\subsection{\texttt{Halo2Verifier.sol} --- the main contract}
\label{sec:halo2-verifier-tpl}

The contract is a thin shell: NatSpec, then includes, then one terminal
assembly block. The \texttt{verifyProof} entry begins with a cheap
ABI-shape guard \emph{before} any generated memory work:
\begin{lstlisting}[style=plain]
function verifyProof(bytes calldata proof, uint256[] calldata instances)
    external view returns (bool) {
    assembly ("memory-safe") {
        if iszero(and(
            eq(calldataload(SELECTOR_BYTES), PROOF_HEAD_OFFSET),
            eq(calldataload(INSTANCES_HEAD_CPTR),
               sub(NUM_INSTANCE_CPTR, SELECTOR_BYTES)))) {
            revert(0, 0)
        }
    }
    // ... VK pin, then the generated Yul pipeline ...
}
\end{lstlisting}
\begin{itemize}
  \item \textcolor{tracegreen}{\textbf{[OK] ABI-shape guard.}} The
        hand-rolled calldata parser is protected by a pre-check that the
        dynamic-argument heads point where the generator expects, so a
        malformed ABI layout reverts before being misread as a proof
        stream.
  \item \textcolor{tracegreen}{\textbf{[OK] Terminal assembly block.}}
        The generated block owns the call-frame memory and is terminal:
        accepted proofs \texttt{return}, all rejects \texttt{revert}.
        Scratch starts at \texttt{0x80}, preserving Solidity's reserved
        prefix and free-memory pointer. (Starting at \texttt{0x80} is
        not sufficient for \texttt{memory-safe} under via-IR; see
        \S\ref{sec:memory-safe}.)
  \item \textcolor{mathpurple}{\textbf{[NOTE] Not a generic verifier.}}
        The NatSpec is explicit that proof layout, VK payload, quotient
        program, and memory layout are all generated for \emph{one}
        \texttt{VerifyingKey}. Application contracts must bind instance
        semantics (state roots, program ids, domain separation)
        separately --- the raw verifier ABI guarantees nothing about them.
\end{itemize}

\subsection{The \texttt{"memory-safe"} annotation --- what it promises and what it relies on}
\label{sec:memory-safe}

Every generated assembly block is tagged \texttt{assembly("memory-safe")}.
This is a \emph{promise to the compiler}: the block does not corrupt
memory the rest of the program needs, so the compiler may trust that and
optimize accordingly (e.g.\ skip defensive re-reads). It does \emph{not}
mean ``the block used Solidity's allocator.''

\paragraph{What the verifier actually does.} Solidity reserves the bottom
128 bytes: scratch space (\texttt{0x00}--\texttt{0x3f}), the
\textbf{free memory pointer} (\texttt{0x40}, normally initialized to
\texttt{0x80}), and the zero slot (\texttt{0x60}); free memory begins at
the free-pointer value and is handed out by bumping \texttt{0x40}. The generated verifier \textbf{bypasses
that allocator entirely}: it uses hard-coded absolute addresses
(\texttt{TRANSCRIPT\_\allowbreak{}MPTR = 0x80}, \texttt{PAIRING\_\allowbreak{}LHS\_\allowbreak{}MPTR =
0x6f20} in B64, \ldots) and \textbf{never advances the free memory pointer}.
(Every \texttt{0x40}/\texttt{0x60} in the emitted code is an
\emph{offset inside a scratch buffer}, e.g.\ \texttt{add(scratch, 0x40)},
not the free pointer itself.)

\paragraph{Why the promise was believed to hold.} Two properties were
taken to make \texttt{memory-safe} true here:
\begin{enumerate}[leftmargin=1.4em,itemsep=1pt]
  \item scratch starts at \texttt{0x80}, so the reserved prefix
        (\texttt{0x00}--\texttt{0x7f}, including the free pointer and
        zero slot) is never touched;
  \item the block is \textbf{terminal} --- it ends in \texttt{return} or
        \texttt{revert}, so no Solidity code runs \emph{after} it.
\end{enumerate}

\paragraph{The load-bearing coupling (the real issue).} The verifier is
doing manual absolute-address memory management that ignores the free
pointer. That is harmless \emph{only because nothing allocates memory
after the block}. The \texttt{memory-safe} tag is a contract that holds
under the terminal-block assumption:

\begin{quote}
If this body were ever inlined into a function that \textbf{keeps running
Solidity after verification}, the verifier's scratch at \texttt{0x80}+
would collide with memory the free pointer hands out next --- and because
the compiler \emph{trusts} the \texttt{memory-safe} tag, it would
\textbf{not} insert the defensive behavior needed to protect against
that collision, yielding silent memory corruption.
\end{quote}

\paragraph{Correction: terminality is not enough under via-IR.}
Solidity's rules for \texttt{memory-safe} blocks hold \emph{inside} the
block too: a block may touch only scratch \texttt{0x00}--\texttt{0x3f},
memory Solidity allocated, and memory at or above the free pointer's value
on entry. The annotation is what lets via-IR's stack-limit evader spill Yul
variables to memory, and it does so here. Compiled with the production
settings (solc 0.8.35, via-IR, 200 runs, Cancun), B64's runtime begins
\texttt{PUSH2 0x08e0 PUSH1 0x40 MSTORE}. The compiler reserved
\texttt{[0x80, 0x8e0)} for spill slots, and that range overlaps the
generator's absolute scratch (the transcript buffer at \texttt{0x80}, the
\texttt{ec\_pairing} scratch at \texttt{0x300}--\texttt{0x600}, \ldots).
Honest proofs verify today only because the live ranges happen to be
disjoint. Nothing checks this (\S\ref{sec:pr-memorysafe}).

In short: the safety does not come from the annotation itself. It comes
from the terminal-block property the annotation is asserting, plus an
unchecked disjointness between compiler spill slots and generator
scratch. The contract's
own NatSpec encodes exactly this warning --- \emph{``Do not inline this
body into Solidity code that continues executing after verification
without reviewing the memory strategy.''}

\begin{itemize}
  \item \textcolor{rustorange}{\textbf{[RISK] Annotation formally
        violated under via-IR.}} \vstatus{CORRECTED} (previously graded
        [OK]). The block writes below the entry free pointer
        (\texttt{0x8e0}) into compiler-owned spill slots. Dropping the tag
        is not an option, because via-IR needs it to evade stack-too-deep.
        \vstatus{OPEN} The fix is to plan the absolute regions above a
        fixed base and assert \texttt{mload(0x40)} $\le$ base on entry, or
        to add a CI check on the compiled free-pointer initializer
        (\S\ref{sec:pr-memorysafe}).
  \item \textcolor{rustorange}{\textbf{[RISK] Annotation is load-bearing
        on terminality.}} The correctness of the tag depends on the
        block never returning to executing Solidity. Any future refactor
        that makes the verifier body non-terminal must either bump the
        free pointer to cover the scratch region or switch the tag to
        \texttt{memory-unsafe} --- otherwise the compiler's trust is
        misplaced. \emph{Audit note:} treat ``terminal'' as a hard
        invariant of the generated body, not an incidental property.
\end{itemize}

\subsection{\texttt{PrecompileSmoke.sol} --- deployment-time fork check}
\label{sec:smoke}

The constructor calls \texttt{require\_\allowbreak{}eip2537\_\allowbreak{}precompiles()}, which
exercises every runtime prerequisite with identity inputs so an
incompatible fork fails \emph{at deployment}, not at first proof:
\begin{lstlisting}[style=plain]
// MCOPY must work (verifier uses it for staging).
mstore(scratch, 0x1234)
mcopy(add(scratch, 32), scratch, 32)
if iszero(eq(mload(add(scratch, 32)), 0x1234)) { revert(0, 0) }
// G1ADD(id,id) -> id, 128-byte return.
staticcall(gas(), G1ADD, scratch, 256, scratch, 128) ...
// Worst-case G1MSM (all identity/zero) -> id, 128-byte return.
staticcall(gas(), G1MSM, msm_scratch, MSM_BYTES, scratch, 128) ...
// PAIRING([(id,id),(id,id)]) -> true, 32-byte return.
staticcall(gas(), PAIRING, scratch, TWO_PAIR_BYTES, scratch, 32) ...
\end{lstlisting}
\begin{itemize}
  \item \textcolor{tracegreen}{\textbf{[OK] Worst-case MSM size.}} The
        G1MSM smoke uses the \emph{largest} MSM input the verifier can
        emit, not a one-pair call --- so a fork that chokes on the real
        input length is caught at deploy.
  \item \textcolor{tracegreen}{\textbf{[OK] Return-shape checks.}} Each
        probe checks \texttt{returndatasize()} and the returned value,
        catching missing precompiles, short returns, and wrong semantics.
  \item \textcolor{codegenblue}{\textbf{[IMPROVE] G2ADD not smoked.}}
        The smoke tests G1ADD/G1MSM/PAIRING but not G2ADD (0x0d). The
        verifier never calls G2ADD (G2 points come from the pinned VK),
        so this is fine --- but a comment noting the deliberate omission
        would prevent an auditor wondering why. \vstatus{OPEN} No such
        comment in \texttt{PrecompileSmoke.sol}; the MCOPY probe shown
        above was added in \texttt{b53e3ebe}.
\end{itemize}

\subsection{\texttt{Constructors.sol} --- dependency pinning}
\label{sec:constructors}

Four constructor variants (embedded/external VK $\times$
embedded/external quotient) all follow the same pattern:
\begin{lstlisting}[style=plain]
require_eip2537_precompiles();
require(authorizedVk.code.length == EXPECTED_VK_LENGTH
     && authorizedVk.codehash == EXPECTED_VK_CODEHASH, "invalid vk");
require(authorizedQuotient.code.length == EXPECTED_QUOTIENT_LENGTH
     && authorizedQuotient.codehash == EXPECTED_QUOTIENT_CODEHASH,
     "invalid quotient");
AUTHORIZED_VK = authorizedVk;
AUTHORIZED_QUOTIENT = authorizedQuotient;
\end{lstlisting}
\begin{itemize}
  \item \textcolor{tracegreen}{\textbf{[OK] Strict codehash+length pin.}}
        Both the VK and the split quotient evaluator are pinned by
        \texttt{codehash} \emph{and} \texttt{length} at deploy, and the
        verifier re-checks before every proof (Sec.~\ref{sec:vkpayload}).
        The split evaluator is treated with the same strictness as the VK
        because it contains generated verifier logic, not a library.
\end{itemize}

\subsection{\texttt{Halo2VerifyingKey.sol} --- INVALID-prefixed payload}
\label{sec:vk-contract}

The VK contract's runtime is \texttt{INVALID || payload}:
\begin{lstlisting}[style=plain]
// byte 0      : INVALID, so the payload cannot be executed
// byte 1..end : generated VK payload copied by Halo2Verifier
let runtime := PAYLOAD_MPTR
let payload := add(runtime, 0x01)
\end{lstlisting}
\begin{itemize}
  \item \textcolor{tracegreen}{\textbf{[OK] INVALID opcode prefix.}}
        Byte 0 is an unconditional \texttt{INVALID} so a direct call to
        the VK contract cannot execute the payload as code. The verifier
        copies the payload via
        \texttt{extcodecopy(vk, VK\_\allowbreak{}MPTR, 0x01, len)} --- skipping the
        INVALID byte --- and references slots by \texttt{VK\_MPTR + i}.
  \item \textcolor{tracegreen}{\textbf{[OK] Quotient program pinned by
        VK codehash.}} The quotient VM's static program lives in this
        pinned runtime, so it is covered by \texttt{EXPECTED\_\allowbreak{}VK\_\allowbreak{}CODEHASH}
        without verifier-side immediates.
\end{itemize}

\subsection{\texttt{Halo2QuotientEvaluator.sol} --- split evaluator}
\label{sec:quotient-evaluator}

In the split render, the verifier \texttt{staticcall}s this contract,
passing its \emph{raw memory frame} as calldata:
\begin{lstlisting}[style=plain]
// calldata[0..FRAME_LEN) == memory[FRAME_BASE..FRAME_BASE+FRAME_LEN)
// Output frame:
//   word 0: QUOTIENT_MAGIC (version guard)
//   word 1: linearization_expected_eval  (= -nu_y(x))
//   word 2..: simple-selector accumulator scalars
\end{lstlisting}
\begin{itemize}
  \item \textcolor{tracegreen}{\textbf{[OK] Magic-guarded frame.}} The
        output carries a generated \texttt{QUOTIENT\_MAGIC} version word
        so a mismatched evaluator version is detected, not silently
        misread.
  \item \textcolor{tracegreen}{\textbf{[OK] Returns $-\nu_y(x)$, not
        $h(x)$.}} The NatSpec is explicit that the evaluator
        reconstructs the $y$-batched numerator and returns the negated
        expected scalar; the $(1-x^n)$ factor stays on the commitment
        side (Def.~\ref{def:neg-numer}).
  \item \textcolor{mathpurple}{\textbf{[NOTE] Frame-copy trust boundary.}}
        The evaluator copies the verifier's memory frame by absolute
        address; the two contracts must agree on the frame layout. This
        coupling is pinned by the codehash pin, but a layout drift in a
        future codegen change would be a silent correctness bug --- the
        magic word is the only runtime guard.
\end{itemize}

\subsection{\texttt{Pcs.yul} and \texttt{TraceReturn.yul}}
\label{sec:pcs-trace-return}

\texttt{Pcs.yul} is a thin emitter that lays down the generated
multi-prepare sub-blocks (grouped for gas-checkpoint attribution) and
traces the pairing inputs. \texttt{TraceReturn.yul} is the terminal
trace + return:
\begin{lstlisting}[style=yul]
// Rust prints the positive numerator; the verifier stores the negated
// expected scalar, so trace id 36 re-negates:
trace_u256(36, addmod(0, sub(r, mload(QUOTIENT_EVAL_MPTR)), r))
// QUOTIENT_MPTR carries scalar metadata in the fused path; zero the
// unused point words before tracing it as a G1 slot:
mstore(add(QUOTIENT_MPTR, 0x40), 0)
mstore(add(QUOTIENT_MPTR, 0x60), 0)
trace_point(24, QUOTIENT_MPTR)
...
mstore(RETURN_MPTR, 1)
return(RETURN_MPTR, 0x20)
\end{lstlisting}
\begin{itemize}
  \item \textcolor{tracegreen}{\textbf{[OK] Stable trace ids.}} The
        terminal trace uses fixed ids matching the Rust
        \texttt{solidity\_\allowbreak{}trace.rs} labels, so fixture tooling can diff
        Rust vs Solidity executions without circuit knowledge
        (Sec.~\ref{sec:rust-trace}).
  \item \textcolor{mathpurple}{\textbf{[NOTE] \texttt{trace\_id 36}
        re-negation.}} The verifier stores $-\nu_y(x)$ but Rust traces
        the positive $\nu_y(x)$; the trace path re-negates to match.
        This is the same negated-numerator subtlety as
        \S\ref{sec:vm-final-neg} --- correct, but a trap if the trace
        consumer assumes the stored value.
\end{itemize}
